\begin{statement}{A.27}{Proved}{Every odd-prime packet on the stable genus-two cover}
Let
\[
 D:\ z^2=(t^2+4)(t^3+3t),\qquad
 Y=D\setminus\{Q_\infty\},
\]
and let $\ell$ be odd. For every original packet section $s_b$
in Part A and every open subgroup $\Delta_0$ over which its
selected lift to $Y$ is defined,
\[
                 P_{D,\ell}^{\,s_b(\Delta_0)}=I_{D,\ell}.
                                                        \tag{A.27.1}
\]
The same assertion holds for all three reflection quotients
of the stable $S_3$ cover, with their selected cusps. In
particular, at odd primes none of the exceptional packet
parameters left for $D$ in A.24 is exceptional for the actual
fixed-subgroup conclusion.

The algebraic input is an unbounded kernel calculation.
Put $L=\mathbb L_k(a,b,c,d)$ over any characteristic-zero
field, and write
\[
 \omega_1=[a,b],\quad \omega_2=[c,d],\quad
 \Omega=\omega_1+\omega_2.
\]
Let $T_i$ act by $\operatorname{ad}(\omega_i)$ on the two
generators of its own component and by zero on the other
component. For $m_1m_2\ne0$ and any $\lambda\in k$, set
\[
             D_\lambda=m_1T_1+m_2T_2
                                +\lambda\operatorname{ad}(\Omega).
                                                        \tag{A.27.2}
\]
Then
\[
 \ker(D_\lambda:L\longrightarrow L)=
 \begin{cases}
   \mathbb L_k(\omega_1,\omega_2),&\lambda=0,\\
   k\Omega,&\lambda\ne0.
 \end{cases}                                           \tag{A.27.3}
\]
The arithmetic proof uses the full separating monodromy
and one of its actual Galois conjugates.
\end{statement}
\begin{proof}
\paragraph{A tensor-algebra normal form.}
First prove (A.27.3). Work in
\[
 A=k\langle a,b,c,d\rangle=A_1*A_2,
 \qquad A_1=k\langle a,b\rangle,
 \quad A_2=k\langle c,d\rangle,
\]
and put $C=k\langle\omega_1,\omega_2\rangle$.
These two quadratic generators freely generate their
associative subalgebra: their leading words $ab$ and $cd$
give distinct leading words for distinct products.

Order words in $a,b$ lexicographically with $a>b$, within
each fixed length. The leading word of $\omega_1$ is $ab$,
which has no proper self-overlap. Let $\mathcal B_1$ be the
set of nonempty words that neither begin nor end with $ab$.
Every word that is not a power of $ab$ has a unique expression
\[
                        (ab)^r w(ab)^s,
 \qquad w\in\mathcal B_1,
\]
obtained by removing its maximal initial and terminal
powers of $ab$. The absence of a self-overlap guarantees
that these removals do not create an additional ambiguity.
Consequently the elements
\[
 \{\omega_1^n:n\ge0\}\ \sqcup\
 \{\omega_1^r w\omega_1^s:
              w\in\mathcal B_1,\ r,s\ge0\}
                                                        \tag{A.27.4}
\]
are a basis of $A_1$: in each finite-dimensional homogeneous
piece their leading words give a triangular matrix with
diagonal entries one. Thus, as a bimodule,
\[
 A_1=k[\omega_1]\ \oplus\!!\bigoplus_{w\in\mathcal B_1}
                          k[\omega_1]w k[\omega_1].
                                                        \tag{A.27.5}
\]
The analogous statement holds for $A_2$, using $cd$ and
$\mathcal B_2$. The displayed summands with $w$ are free
bimodules; no finiteness assertion for $\mathcal B_i$ is
being made.

Apply these two bases to the usual alternating-block
basis of the free product $A_1*A_2$. A basis word either
belongs to $C$, or has a first nonboundary core
$w\in\mathcal B_i$. All the pure $\omega$-blocks preceding
that core, together with the left power of $\omega_i$
in its own block, form an arbitrary word $u$ in $C$.
What follows is a tail beginning with that fixed core.
This is a unique normal form. In particular there are
tail spaces $M_i$ and a vector-space direct sum
\[
                    A=C\ \oplus\ (C\otimes M_1)
                           \ \oplus\ (C\otimes M_2).
                                                        \tag{A.27.6}
\]
Here a tail begins with a nonboundary core of type $i$,
and the tensor notation records its uniquely determined
prefix, not an assertion that $A$ is a free algebra over $C$.

Give a prefix its length as a word in the two letters
$\omega_1,\omega_2$. The derivation
$D_0=m_1T_1+m_2T_2$ kills every pure boundary block.
On the first core of type $i$ it is
\[
                      D_0(w)=m_i(\omega_iw-w\omega_i).
\]
The first term increases the prefix length by one, by
right multiplication of the prefix by $\omega_i$.
The second term and the action on all later cores preserve
that prefix length. These operations cannot remove or
change the first core: multiplying by $\omega_i$ changes
only the corresponding left or right power in (A.27.4).
Likewise the left part of $\lambda\operatorname{ad}(\Omega)$
increases prefix length by one, whereas its right part
preserves it. Appending a pure boundary word does not
remove a nonboundary core.

It follows that the part of $D_\lambda$ increasing the
prefix length on $C\otimes M_i$ is
\[
           (\lambda L_\Omega+m_i R_{\omega_i})\otimes1.
                                                        \tag{A.27.7}
\]
This operator is injective on $C$. If $\lambda=0$, this
is just nonzero right multiplication in a free algebra.
If $\lambda\ne0$, let $j\ne i$ and take the largest
number of occurrences of $\omega_j$ in a nonzero finite
linear combination of words. In its image under (A.27.7),
the terms with one more such occurrence come only from
$\lambda L_{\omega_j}$, and cannot cancel. This also
proves injectivity after tensoring with any tail space.

Now choose the largest prefix length in a nonzero element
of $C\otimes M_i$. Its image of prefix length one greater
is nonzero by (A.27.7); terms of smaller prefix length
cannot contribute there. Thus no nonzero element involving
a nonboundary core is killed by $D_\lambda$, and
\[
                         \ker(D_\lambda|_A)\subset C.
                                                        \tag{A.27.8}
\]
On $C$ the derivation is simply
$\lambda\operatorname{ad}(\Omega)$. For $\lambda=0$ its
kernel is all of $C$. For $\lambda\ne0$ its kernel is
$k[\Omega]$: replace the two free generators of $C$ by
$\Omega,\omega_1$, and use the elementary centralizer
calculation for a free generator. Finally $\omega_1$ and
$\omega_2$ are primitive in the tensor Hopf algebra.
The primitive elements of $C$ form
$\mathbb L_k(\omega_1,\omega_2)$, and those of
$k[\Omega]$ form $k\Omega$. Intersecting (A.27.8) with
$L$ proves (A.27.3) in every degree.

\paragraph{The cusp lies on the rational bridge.}
Use the stable reduction of A.25 at a place above $2$.
Its two elliptic components $E_1,E_2$ have degrees $1,2$
under the extension of
\[
 f:D\longrightarrow E,\qquad
 x_E=-i(t^3+3t)/2,\quad
 y_E=-(1+i)z(t^2+1)/4,
\]
and $f^*(O)=3Q_\infty$ on the generic fibre. Neither
component is contracted, so the Cartier divisor $f^*(O)$
on the normal stable total space has no vertical component.
Its Weil divisor is therefore $3\overline{Q_\infty}$.
If the cusp section met a smooth point of just one
elliptic component, the support of this divisor would
miss the other component entirely. The restriction of
$f$ to that other elliptic component is nonconstant and
surjective, so it must have a preimage of $O$, a
contradiction. Hence the cusp specializes to the node.

The stable pointed model consequently inserts a rational
component carrying the cusp between $E_1$ and $E_2$.
After a further finite extension and regular subdivision,
the distances from this rational component to the two
elliptic components are both strictly positive. Denote
their total edge lengths by $e_1,e_2>0$. Choose an ordinary
base point reducing to a smooth point of the rational
component away from its two nodes and the cusp. Such a
point exists after a finite unramified extension and lifts
by smoothness.

For this affine curve and ordinary base point,
\cite[Definition 3.1.5 and Theorem 3.1.6]{fghBettsDogra}
gives the based edge-twist description. The rational root
component has three boundary loops: the cusp and the two
nodes. Its vertex group is fixed by monodromy. Transport
the two elliptic handle pairs to that root. The first
nonzero ordinary Johnson derivation of tame inertia is
\[
                            e_1T_1+e_2T_2,
                                                        \tag{A.27.9}
\]
with a possible simultaneous sign depending on the
chosen inertia convention. Each $T_i$ fixes both
$\omega_1,\omega_2$, and hence the cusp class $\Omega$.
There is no first-degree Johnson term: every conjugating
boundary word starts in lower-central degree two.

\paragraph{Actual packet monodromy and its only inner term.}
Choose a finite number field $L$ defining all these maps,
cusps and CM isogenies, such that $\Delta\cap G_L$ is
contained in the prescribed $\Delta_0$. Enlarge $L$ to
realize the local semistable situation above, and put
\[
                      \Gamma=\overline{[G_L,G_L]}.
\]
This is a normal subgroup of $G_L$ contained in
$\Delta_0$. Its actual packet action on $P_{D,\ell}$ is
trivial on homology, since $J_D\sim E^2$ and the CM
characters are abelian. By the tangential hyperelliptic
equivariance established in A.24, the first Johnson
degree of this action is zero: the involution acts as
$-1$ on homology and $+1$ on its second Lie degree.
The square of its based lift may be inner; only the
induced graded signs are used here.

Let $F$ be a local Frobenius above $2$ and $\tau$ a tame
$\ell$-inertia element with nonzero monodromy. Then
$\delta=[F,\tau]\in\Gamma$. Its affine outer monodromy
is $(q-1)$ times that of $\tau$, by the tame relation,
where $q>1$ is the residue cardinality. In the semistable
prime-to-two unipotent outer action wild inertia is
trivial. Thus an outer representative for the first
nonzero derivation of $\delta$ is
\[
             m_1T_1+m_2T_2,
                \qquad m_i=(q-1)e_i\ne0,
                                                        \tag{A.27.10}
\]
up to the same nonzero normalization for both coefficients.

Compare this ordinary-base-point representative with
the actual selected packet action. They have the same
\emph{affine} outer action, so their difference is inner.
Their first Johnson degrees are both zero. Consequently
the linear term of an inner conjugator is zero: a linear
free-Lie element commuting with all four generators is
zero. In degree two the difference therefore has the form
$\operatorname{ad}(w_2)$, with $w_2\in L_2$.
The actual packet fixes its actual cusp element on
$\Gamma\subset\Delta$; since its first Johnson degree
vanishes, its degree-two derivation kills the leading
cusp class $\Omega$. The representative (A.27.10) kills
$\Omega$ as well. Hence
\[
                              [w_2,\Omega]=0.
\]
The free-Lie centralizer calculation proved in A.24
implies $w_2=\lambda\Omega$. We have obtained the
\emph{actual based} degree-two derivation (A.27.2), with
both $m_i\ne0$ and an arbitrary $\lambda\in\Q_\ell$.
This comparison requires no tangential extension of the
ordinary-base-point theorem of Betts--Dogra.

\paragraph{A Galois conjugate removes the case $\lambda=0$.}
If $\lambda\ne0$, (A.27.3) already gives the required
graded kernel. Suppose $\lambda=0$. Since $\Gamma$ is
normal in $G_L$, the degree-two image of its actual
action is stable under conjugation by $s_b(G_L)$.
Conjugation on this first nonzero degree depends only
on the action on homology. We exhibit one such
conjugate whose kernel meets that of $D_0$ only in
$\Q_\ell\Omega$.

Extend coefficients to split the CM field. Choose
symplectic local component bases $a,b$ and $c,d$ with
\[
              f_*a=x,\quad f_*b=y,\quad
              f_*c=x,\quad f_*d=2y,
                                                        \tag{A.27.11}
\]
where $x,y$ is a CM eigenbasis on $E$. These
normalizations respect the polarization because the
two component maps have degrees one and two.
Adjointness identifies the original elliptic subspace
with $\langle a+2c,b+d\rangle$, and the complementary
subspace with $\langle a-c,2b-d\rangle$.
The automorphism $\rho(t,z)=(-t,-iz)$ satisfies
$f\rho=[-i]f$. It acts by $[-i]$ on the original
elliptic factor and by $[i]$ on the complementary one,
as seen on $dt/z$ and $t\,dt/z$. It preserves the two
special components because their map degrees differ.
Thus a symplectic CM eigenbasis is
\[
 p_1=a+2c,\quad p_2=(2b-d)/3,\qquad
 q_1=(b+d)/3,\quad q_2=c-a,
                                                        \tag{A.27.12}
\]
where the $p$'s have one CM character $\psi$ and the
$q$'s the other $\bar\psi$. These are rational
coordinates, also when $\ell=3$; no integral division
by three is asserted.

Put $P=[p_1,p_2]$, $Q=[q_1,q_2]$,
$A=[p_1,q_1]$, and $B=[p_2,q_2]$. Direct substitution
gives
\[
 \Omega=A+B,\qquad
 \omega_1=\tfrac13(P+A+2B+2Q),\qquad
 \omega_2=\tfrac13(-P+2A+B-2Q).
                                                        \tag{A.27.13}
\]
Take an actual Frobenius $g\in G_L$ at an ordinary
good prime of $E$ split in its CM field and different
from $\ell$. Its two CM eigenvalues $\alpha,\bar\alpha$
have ratio of infinite order: as algebraic numbers,
their prime-ideal supports above that rational prime
are disjoint. Passing to a positive residue power
does not change this fact. If $t$ is the action of
$g$ on homology, (A.27.13) shows
\[
 U\cap tU=k\Omega,\qquad
                         U=\langle\omega_1,\omega_2\rangle.
                                                        \tag{A.27.14}
\]
Indeed both planes contain $\Omega$ up to its
cyclotomic scalar, whereas the ratio of their $P$ and
$Q$ coefficients is changed by
$(\alpha/\bar\alpha)^2\ne1$.

Inside $V^{\otimes 2r}$, group the tensor factors in
pairs. For subspaces $U,U'\subset V^{\otimes2}$,
\[
                  U^{\otimes r}\cap(U')^{\otimes r}
                            =(U\cap U')^{\otimes r}.
                                                        \tag{A.27.15}
\]
This follows by choosing a vector-space basis adapted
to $U\cap U'$, then to $U$ and $U'$, and finally to
their sum. Hence the associative algebras freely
generated by $U$ and $tU$ intersect in $k[\Omega]$.
Their primitive parts intersect in $k\Omega$.
By (A.27.3), the two actual degree-two derivations
from $\delta$ and $g\delta g^{-1}$ therefore have
common free-Lie kernel exactly $k\Omega$.
Descent from the coefficient extension gives the
same assertion over $\Q_\ell$.

\paragraph{From the graded obstruction to actual fixed elements.}
Let $u\in P_{D,\ell}$ be fixed by $s_b(\Delta_0)$.
The closed lower-central filtration of this free
pro-$\ell$ group is separated and has torsion-free
free-Lie graded pieces over $\Z_\ell$. If $u\ne1$,
its first nonzero rational graded class is killed
by the actual degree-two derivations above. It is
therefore $a\Omega$ in degree two. The integral
leading class of the selected cusp element $c_D$
is primitive: in an integral symplectic basis it
is $[a,b]+[c,d]$, up to orientation. Thus $a\in\Z_\ell$.

The actual cusp power $c_D^a$ is fixed by
$s_b(\Delta_0)$. Multiplying $u$ by its inverse
removes that leading term and leaves a fixed
element in $\gamma_3P_{D,\ell}$. A nontrivial
remainder would have a leading class of degree
at least three in the same common kernel, which
is impossible. Hence $u=c_D^a$. The reverse
inclusion in (A.27.1) follows from the trivial
cyclotomic action on $\Delta_0$.

The three reflection quotients and their chosen
cusp lifts are conjugate after a finite extension,
so the argument applies to each. The proof uses
the place above two and therefore asserts
(A.27.1) only for odd $\ell$.
\end{proof}
